\section{Experimental Design}

The evaluation rests on three studies that approach the same question from
complementary methodological directions.

\textbf{Experiment~1} is observational. Four review-generation modes are run
over 35 pull requests and scored against a 25-criterion rubric by a panel
of five LLM judges, with the majority verdict as the cell of record. A
language-model panel is used because the rubric requires 3\,500 binary
judgements per scorer (140 reviews $\times$ 25 criteria); a human gold
standard at that scale was not feasible, so human judgement is instead
reserved for the primary criterion in the third study. Experiment~1's
strength is direct evidence on pull-request reviews; its weakness is that it
establishes no causal mechanism and that the rubric measures coverage rather
than correctness.

\textbf{Experiment~2} is controlled. A defect is injected so that its
consequences lie in a different file from the edit; the outcome is whether the
review names a genuinely broken dependent site. Here the ground truth is known,
so a causal claim becomes available. The cost is that the defects are synthetic.

\textbf{The human study} asks people. Twenty-seven raters compare blinded
baseline and graph-augmented reviews on six pull requests, scoring
which review better identifies affected code and which is more useful overall.
Its strength is independent human judgement; its weakness is a smaller
stimulus set,
whose graph edges are resolved from import statements rather than from the
Joern CPG.

Each design covers a weakness of the others, although none removes them
all. The headline arms of Experiments~1 and~2 share their generator, judge
panel, and aggregation rule so that differences cannot be attributed to
those choices. The human study corroborates the panel's primary finding
against independent human judgement on a separate stimulus set; it does not
validate the panel itself.

\subsection{Research Questions}
\label{sec:rqs-addressed}

Experiment~1 addresses RQ1 by comparing rubric coverage across the four review
modes. Experiment~2 addresses RQ2 by testing detection of consequences outside
the diff; its component analyses explain which forms of structural evidence
support that capability. The human study addresses RQ3 through blinded
pairwise comparisons.

\subsection{Dataset}
\label{sec:dataset}

The evaluation uses 35 pull requests drawn from the pool curated by
Mariotto et~al.~\cite{mariotto2025esem}, sampling from five repositories:
\texttt{grafana/grafana}, \texttt{apache/kafka},
\texttt{scikit-learn/scikit-learn}, \texttt{godotengine/godot} and
\texttt{jenkinsci/jenkins}, spanning Go, TypeScript, Java, Scala, Python, and
C++.\footnote{Five Go-only pull requests from \texttt{grafana/grafana} are
excluded because Joern's Go frontend does not emit method-level nodes,
preventing construction of the call-graph edges required by the \kg{} arm.}
Table~\ref{tab:dataset-composition} gives the composition.

\begin{table}[htbp]
    \centering
    \caption{Composition of the evaluation dataset.}
    \label{tab:dataset-composition}
    \begin{tabular}{lrl}
        \toprule
        Repository & Total & Languages \\
        \midrule
        \texttt{grafana/grafana}             &  8 & Go, TS \\
        \texttt{apache/kafka}                &  8 & Java, Scala \\
        \texttt{scikit-learn/scikit-learn}   &  8 & Python \\
        \texttt{godotengine/godot}           &  6 & C++ \\
        \texttt{jenkinsci/jenkins}           &  5 & Java \\
        \midrule
        Total                                & 35 & \\
        \bottomrule
    \end{tabular}
\end{table}

Every candidate was screened against seven criteria, all properties of the pull
request rather than of any model output. A pull request is included only if it
was merged, is not a revert or backport, changes code (not only documentation),
touches only languages the parser supports, has a diff under 50\,000 characters, carries a
description of at least 100~characters, and does not have a purely cosmetic
title. This guarantees that all four modes, including \baseline{}, receive the
maintainer's own description.

At $n=35$, a two-sided paired permutation test at $\alpha = 0.05$ reaches
80\,\% power at
$d_z \approx 0.47$ and 70\,\% at $d_z \approx 0.42$,
bracketing the effects reported in Chapter~\ref{chap:results}.

\subsection{Rubric}
\label{sec:rubric}

Each review is scored against 25 binary criteria grouped into eight categories
(Table~\ref{tab:rubric-categories}). A criterion is satisfied if the review
addresses it; the score is the count of criteria satisfied.

\begin{table}[htbp]
    \centering
    \caption{The eight rubric categories. The KG-relevant subset comprises the
             nine criteria structural data can in principle inform.}
    \label{tab:rubric-categories}
    \begin{tabular}{llrl}
        \toprule
        Category & Criteria & $n$ & KG-relevant \\
        \midrule
        Functionality \& Integration & F1--F4 & 4 & F3, F4 \\
        Tests \& Verification        & T1--T3 & 3 & T1, T2, T3 \\
        Readability \& Structure     & R1--R3 & 3 & --- \\
        Maintainability \& Design    & M1--M3 & 3 & M1, M3 \\
        Consistency \& Style         & C1--C2 & 2 & C2 \\
        Performance                  & P1--P2 & 2 & --- \\
        Security \& Robustness       & S1--S3 & 3 & --- \\
        Review Quality (meta)        & Q1--Q5 & 5 & Q2 \\
        \midrule
        Total                        &        & 25 & 9 \\
        \bottomrule
    \end{tabular}
\end{table}

Every category is grounded in published sources: empirical studies of
review~\cite{bacchelli2013,bosu2015,sadowski2018}, evaluation frameworks for
generated comments~\cite{li2022codereviewer,hong2024deepcrceval}, and
ISO/IEC~25010~\cite{iso25010}.

\subsubsection{The KG-Relevant Subset}

Nine criteria are designated KG-relevant. The designation was fixed before
data collection and appears in the project's analysis plan committed to the
repository. A criterion qualifies if structural repository data can
substantively inform an answer to it. F4, for instance, asks whether the
review warns about impact on dependent code, which the caller fields can
answer, whereas R1 (naming and clarity) they cannot. This yields a testable
prediction: if the knowledge graph works as intended, its effect should
concentrate in these nine criteria. An independent re-annotation achieved
$\kappa = 0.615$ (84\,\% raw agreement).

\subsubsection{Rubric Discrimination}

To confirm the rubric distinguishes good from bad reviews, five deliberately
degraded review styles were scored on five pull requests. Four scored 5--9
points below the real \baseline{}. The fifth, a hallucinated review citing
fabricated files, scored 9.2 against the real \baseline{}'s 9.0. The rubric
measures whether a review \emph{addresses} a concern, not whether what it says
is \emph{true}. This is why every result is framed as coverage, not
correctness, and why Experiment~2's oracle is needed.

An operational note: seven of the 25 criteria sit at $0\,\%$ or $100\,\%$ in
both arms (Appendix, Table~\ref{tab:exp1-criteria}) and therefore cannot
register a difference in either direction, diluting the aggregate. This is
one reason the nine-criterion subscale is reported alongside the total.

\subsection{Experiment 2: Controlled Defect Injection}
\label{sec:injection-design}

Experiment~2 was designed to close the gap between coverage and correctness on
a narrow question: when a defect's consequences lie outside the edited file,
does context let the reviewer name them?

\subsubsection{Defect Taxonomy}

Injections are stratified into five structural bands (consequence in another
file) and two local control bands (consequence in the same file).

\begin{table}[htbp]
    \centering
    \caption{Injection bands. Structural bands require cross-file reasoning;
             control bands serve as negative controls, where no graph
             advantage is expected.}
    \label{tab:injection-bands}
    \small
    \begin{tabular}{llll}
        \toprule
        Band & Operator & Breaks at & Relationship \\
        \midrule
        S1 & Symbol rename          & Importers          & imports \\
        S2 & Signature change       & Callers            & calls \\
        S3 & Return-type change     & Consumers          & calls, data flow \\
        S4 & Base-class change      & Subclasses         & inheritance \\
        S5 & Import removal         & Importers          & imports \\
        \midrule
        L1 & Off-by-one             & Same file          & --- \\
        L2 & Removed null check     & Same file          & --- \\
        \bottomrule
    \end{tabular}
\end{table}

Band sizes are uneven (S1 9, S2 9, S3 1, S4 6, S5 3; six in each control
band), so S3 and S5 support no band-level claims.

The control bands carry the argument: a knowledge graph that improved detection
everywhere would be evidence of a context-volume effect rather than structural
reasoning. The prediction is asymmetric: advantage on structural bands,
parity on local. This is what makes the result falsifiable.

Forty injections were realised: 28 structural and 12 control, across three
of the Experiment~1 repositories: \texttt{scikit-learn/scikit-learn}
(Python), \texttt{apache/kafka} (Java) and \texttt{grafana/grafana}
(TypeScript). The four modes of Experiment~1 are carried over with the same
prompts, builder and renderer (Table~\ref{tab:config}). Two further arms are
included for the mechanism argument: an \emph{idealised} graph that is given
the oracle's own list of true dependents, and a graph with inheritance edges
added (\kg{} + inheritance).

Detection is binary and adjudicated in two stages. First, the review text is
searched for exact full-path or basename matches against the statically
resolved true dependents; if none is found, the verdict is deterministically
\emph{false} and no model call is made. When at least one true dependent is
mentioned, the review passages and matched evidence are sent to each judge,
which decides only whether the review links the mentioned file to the injected
change as broken or requiring an update. The majority of five judges is the
cell of record, with ties resolved to \emph{false}.

This two-stage procedure was finalised after an audit of the initial judge
prompt, which required judges to both locate filenames and assess causality in
a single step. Retained rationales exposed false negatives: reviews that
named a true dependent were sometimes judged as not detecting it. The
correction is post-hoc; the original judgments are preserved in the audit
trail and both sets of verdicts are available for recomputation.

\subsubsection{Component Analyses}

Two pre-registered follow-ups isolate the structural evidence supplied to the
reviewer. On the 28 structural injections, one ablation arm receives only the
dependent-file list (without edge annotations) and another only the resolved
call edges; both use the final five-judge adjudication. Two further
component probes, a related-test band and a context-relevance control, use
the earlier three-judge panel and are reported in
Appendix~\ref{app:supporting-components}.

\subsection{Statistical Analysis}
\label{sec:stats}

For Experiment~1's rubric scores, confidence intervals are percentile
bootstrap intervals ($B = 10\,000$).
Significance is assessed by paired sign-flip permutation tests
($B = 20\,000$), two-sided, reported alongside $d_z$. The pairing is by pull
request.

The governing RQ1 test is the pre-specified \kg{}--\baseline{} contrast on the
nine KG-relevant criteria. A rank-based multiplicity sensitivity and
per-criterion corrections are retained in
Appendix~\ref{app:supporting-multiplicity}, rather than treated as additional
headline tests.

For Experiment~2's binary detection outcomes, exact McNemar tests compare
paired arms and Wilson intervals replace bootstrap intervals (the appropriate
method for binomial proportions at small $n$).
Decision rules were fixed before the run: $+0.30$ detection
rate advantage at $p < 0.05$ for \kg{} on structural bands, parity on control
bands within $[-0.15, +0.15]$, and $+0.30$ for \kg{} + inheritance on the
structural bands.
